\documentclass[11pt]{article}
\usepackage[margin=1in]{geometry}
\usepackage{amsmath,amssymb}
\usepackage{enumitem}
\usepackage[T1]{fontenc}

\title{COMP2300 Set 5 Practice Exam}
\author{Topics: Multi-Cycle Microarchitecture, FSM Control, Performance,\\ Microprogramming, Pipelining Intro}
\date{}

\begin{document}
\maketitle

\noindent\textbf{Time:} 120 minutes \hfill
\textbf{Total:} 100 marks

\vspace{0.5em}
\noindent\textbf{Instructions}
\begin{itemize}[leftmargin=1.5em]
  \item Answer in English.
  \item Part A is multiple-choice (concept-focused).
  \item Part B is computational. Show key steps.
  \item Part C contains past exam questions with sources.
  \item Assume 32-bit ARM unless stated otherwise.
\end{itemize}

\section*{Part A: Multiple Choice (30 marks, 3 marks each)}
Choose exactly one option for each question.

\begin{enumerate}[label=\textbf{A\arabic*.}]
  \item The main motivation for moving from a single-cycle CPU to a multi-cycle CPU is to:
  \begin{enumerate}[label=(\Alph*)]
    \item Guarantee CPI $= 1$ for every instruction
    \item Shorten the clock period by breaking execution into states
    \item Remove the need for a control unit
    \item Make instruction encoding simpler
  \end{enumerate}

  \item From the ISA point of view, the intermediate states used by a multi-cycle implementation are:
  \begin{enumerate}[label=(\Alph*)]
    \item Architecturally visible to the assembly programmer
    \item Required to be stored in general-purpose registers
    \item Programmer-invisible microarchitectural state
    \item Part of the instruction opcode
  \end{enumerate}

  \item Which of the following is a downside of a multi-cycle design?
  \begin{enumerate}[label=(\Alph*)]
    \item Hardware resources cannot be reused
    \item Additional state registers and sequencing overhead are needed
    \item Branches cannot be implemented
    \item CPI is fixed at 1
  \end{enumerate}

  \item In the final multi-cycle ARM controller discussed in Set 5, which instruction class takes 5 cycles?
  \begin{enumerate}[label=(\Alph*)]
    \item \texttt{B}
    \item Data-processing
    \item \texttt{STR}
    \item \texttt{LDR}
  \end{enumerate}

  \item When \texttt{R15} is read as a source operand, the register file should provide:
  \begin{enumerate}[label=(\Alph*)]
    \item The current PC value
    \item PC + 4
    \item PC + 8
    \item The branch target address
  \end{enumerate}

  \item For an ARM \texttt{B} instruction, \texttt{ExtImm} is formed by:
  \begin{enumerate}[label=(\Alph*)]
    \item Zero-extending \texttt{imm12}
    \item Sign-extending \texttt{imm24} after shifting left by 2
    \item Sign-extending \texttt{imm8}
    \item Adding 8 directly to the PC
  \end{enumerate}

  \item In microprogrammed control, the control signals for the current state are stored in a:
  \begin{enumerate}[label=(\Alph*)]
    \item Branch predictor
    \item Microinstruction
    \item Register file entry
    \item Cache line
  \end{enumerate}

  \item The component that determines the address of the next microinstruction is the:
  \begin{enumerate}[label=(\Alph*)]
    \item Microsequencer
    \item ALU
    \item Data memory
    \item Extend block
  \end{enumerate}

  \item The main limitation of a multi-cycle processor that motivates pipelining is:
  \begin{enumerate}[label=(\Alph*)]
    \item Instructions become variable-width
    \item Hardware resources are often idle, limiting concurrency
    \item Immediate values become too small
    \item PC can no longer be updated on branches
  \end{enumerate}

  \item Pipelining primarily improves:
  \begin{enumerate}[label=(\Alph*)]
    \item The ISA semantics
    \item Throughput by overlapping different stages of different instructions
    \item The width of immediate fields
    \item The number of architectural registers
  \end{enumerate}
\end{enumerate}

\section*{Part B: Computational Problems (50 marks)}
\begin{enumerate}[label=\textbf{B\arabic*.}]
  \item \textbf{State sequences and cycle counts (8 marks)}\\
  In the final multi-cycle FSM, the available states include \texttt{Fetch}, \texttt{Decode}, \texttt{MemAdr}, \texttt{MemRead}, \texttt{MemWB}, \texttt{MemWrite}, \texttt{ExecuteI}, \texttt{ExecuteR}, \texttt{ALUWB}, and \texttt{Branch}.
  \begin{enumerate}[label=(\alph*)]
    \item Write the state sequence and cycle count for each of the following instruction classes: \texttt{LDR}, \texttt{STR}, data-processing immediate, and \texttt{B}.
    \item How many cycles are needed to execute the following 6-instruction sequence on this processor? What is the average CPI?
\begin{verbatim}
LDR R0, [R1, #0]
ADD R2, R2, #1
STR R2, [R1, #4]
B label
ADD R3, R3, R4
LDR R5, [R6, #8]
\end{verbatim}
  \end{enumerate}

  \item \textbf{PC, R15, and branch target calculation (8 marks)}\\
  An instruction is located at address \texttt{0x80A0}.
  \begin{enumerate}[label=(\alph*)]
    \item If the instruction is \texttt{ADD R1, R15, R2}, what value should be supplied when \texttt{R15} is read in the register-read step?
    \item If the instruction is \texttt{B} with \texttt{imm24 = 0x000003}, compute \texttt{ExtImm} and the branch target address.
  \end{enumerate}

  \item \textbf{Decoder outputs (8 marks)}\\
  The decoder uses the following relations:
  \[
    \texttt{RegSrc0 = (Op == 10\_2)}, \qquad
    \texttt{RegSrc1 = (Op == 01\_2)}, \qquad
    \texttt{ImmSrc = Op}
  \]
  Fill in the outputs for each instruction class below:
  \begin{center}
  \begin{tabular}{c|c|c|c}
  Instruction class & \texttt{Op} & \texttt{RegSrc1:0} & \texttt{ImmSrc} \\ \hline
  \texttt{LDR} & 01 & ? & ? \\
  \texttt{STR} & 01 & ? & ? \\
  DP immediate & 00 & ? & ? \\
  DP register & 00 & ? & ? \\
  \texttt{B} & 10 & ? & ? \\
  \end{tabular}
  \end{center}

  \item \textbf{Average CPI and total cycles (8 marks)}\\
  A workload contains 25\% loads, 10\% stores, 13\% branches, and 52\% data-processing instructions.
  Assume:
  \[
    CPI_{LDR}=5,\quad CPI_{STR}=4,\quad CPI_B=3,\quad CPI_{DP}=4
  \]
  \begin{enumerate}[label=(\alph*)]
    \item Compute the average CPI.
    \item If the program has 50 million instructions, how many total cycles does it take?
  \end{enumerate}

  \item \textbf{Critical path and clock frequency (8 marks)}\\
  Assume the multi-cycle critical-path formula is
  \[
    T_{c2} = t_{pcq} + 2t_{mux} + \max(t_{ALU} + t_{mux}, t_{mem}) + t_{setup}
  \]
  and the delays are:
  \[
    t_{pcq}=40\text{ ps},\quad t_{mux}=25\text{ ps},\quad
    t_{ALU}=120\text{ ps},\quad t_{mem}=200\text{ ps},\quad
    t_{setup}=50\text{ ps}
  \]
  \begin{enumerate}[label=(\alph*)]
    \item Compute \(T_{c2}\).
    \item Compute the corresponding maximum clock frequency in GHz.
  \end{enumerate}

  \item \textbf{Execution time (10 marks)}\\
  A program executes \(100 \times 10^9\) instructions on a multi-cycle ARM processor with:
  \[
    CPI = 4.12,\qquad T_c = 340\text{ ps}
  \]
  \begin{enumerate}[label=(\alph*)]
    \item Compute the execution time in seconds.
    \item Suppose a revised design keeps the same CPI but reduces the clock period to \(300\) ps. What is the new execution time?
    \item What is the speedup of the revised design over the original design?
  \end{enumerate}
\end{enumerate}

\section*{Part C: Past Exam Questions (Source-Tagged) (20 marks)}
\begin{enumerate}[label=\textbf{C\arabic*.}]
  \item \textbf{ISA vs Microarchitecture (Source: exam24px1.pdf, Q1f--Q1i, adapted) (5 marks)}\\
  For each item below, classify it as \textbf{ISA} or \textbf{Microarchitecture}.
  \begin{enumerate}[label=(\alph*)]
    \item Whether the ALU uses a ripple-carry adder or a carry-save adder
    \item Instruction width
    \item Whether the CPU uses branch prediction
    \item Whether the control unit is hardwired or microprogrammed
  \end{enumerate}

  \item \textbf{Multi-Cycle Execution Time (Source: exam2025-main.pdf, Q12) (5 marks)}\\
  A program consists of 2 billion instructions, of which 50\% are ALU instructions, 25\% are memory instructions, and 25\% are branch instructions.
  On a 1 GHz multi-cycle CPU, an ALU instruction takes 6 cycles, each memory instruction takes 8 cycles, and each branch instruction takes 4 cycles.
  What is the execution time in seconds?

  \item \textbf{Single-Cycle vs Pipelining (Source: exam24px2.pdf, Q6, adapted) (5 marks)}\\
  State two drawbacks of a single-cycle processor microarchitecture. Explain how instruction pipelining addresses these drawbacks. Also define instruction-level parallelism (ILP), and name two types of microarchitectures that exploit ILP.

  \item \textbf{ILP Ranking (Source: 2023\_Final\_Exam.pdf, Bonus Q1, adapted) (5 marks)}\\
  Rank the following microarchitectures from \textbf{highest} to \textbf{lowest} potential to exploit instruction-level parallelism. Briefly justify your ranking.
  \begin{enumerate}[label=(\Alph*)]
    \item Scalar in-order pipeline with forwarding and hazard detection
    \item Scalar in-order pipeline without forwarding or hazard detection (compiler inserts NOPs)
    \item Multi-cycle processor
  \end{enumerate}
\end{enumerate}

\end{document}
